\section{Flat modules and flat ring maps}
\label{section-flat}

\noindent
One often used result is that if $M = \colim_{i\in \mathcal{I}} M_i$
is a colimit of $R$-modules and if $N$ is an $R$-module then
$$
M \otimes N
=
\colim_{i\in \mathcal{I}} M_i \otimes_R N,
$$
see Lemma \ref{lemma-tensor-products-commute-with-limits}.
This property is usually expressed by saying
that {\it $\otimes$ commutes with colimits}.
Another often used result is that if $0 \to N_1 \to N_2 \to N_3 \to 0$
is an exact sequence and if $M$ is any $R$-module, then
$$
M \otimes_R N_1
\to
M \otimes_R N_2
\to
M \otimes_R N_3
\to
0
$$
is still exact, see Lemma \ref{lemma-tensor-product-exact}.
Both of these properties tell us that the functor
$N \mapsto M \otimes_R N$ {\it is right exact}.
See Categories, Section \ref{categories-section-exact-functor}
and Homology, Section \ref{homology-section-functors}.
An $R$-module $M$ is flat if $N \mapsto N \otimes_R M$ is also left exact,
i.e., if it is exact. Here is the precise definition.

\begin{definition}
\label{definition-flat}
Let $R$ be a ring.
\begin{enumerate}
\item An $R$-module $M$ is called {\it flat} if whenever
$N_1 \to N_2 \to N_3$ is an exact sequence of $R$-modules
the sequence $M \otimes_R N_1 \to M \otimes_R N_2 \to M \otimes_R N_3$
is exact as well.
\item An $R$-module $M$ is called {\it faithfully flat} if the
complex of $R$-modules
$N_1 \to N_2 \to N_3$ is exact if and only if
the sequence $M \otimes_R N_1 \to M \otimes_R N_2 \to M \otimes_R N_3$
is exact.
\item A ring map $R \to S$ is called {\it flat} if
$S$ is flat as an $R$-module.
\item A ring map $R \to S$ is called {\it faithfully flat} if
$S$ is faithfully flat as an $R$-module.
\end{enumerate}
\end{definition}

\noindent
Here is an example of how you can use the flatness condition.

\begin{lemma}
\label{lemma-flat-intersect-ideals}
Let $R$ be a ring. Let $I, J \subset R$ be ideals. Let $M$ be a flat
$R$-module. Then $IM \cap JM = (I \cap J)M$.
\end{lemma}

\begin{proof}
Consider the exact sequence $0 \to I \cap J \to R \to R/I \oplus R/J$.
Tensoring with the flat module $M$ we obtain an exact sequence
$$
0 \to (I \cap J) \otimes_R M \to M \to M/IM \oplus M/JM
$$
Since the kernel of $M \to M/IM \oplus M/JM$ is equal to
$IM \cap JM$ we conclude.
\end{proof}

\begin{lemma}
\label{lemma-colimit-flat}
Let $R$ be a ring. Let $\{M_i, \varphi_{ii'}\}$ be a directed system of
flat $R$-modules. Then $\colim_i M_i$ is a flat $R$-module.
\end{lemma}

\begin{proof}
This follows as $\otimes$ commutes with colimits and because
directed colimits are exact, see
Lemma \ref{lemma-directed-colimit-exact}.
\end{proof}

\begin{lemma}
\label{lemma-composition-flat}
A composition of (faithfully) flat ring maps is
(faithfully) flat.
If $R \to R'$ is (faithfully) flat, and $M'$ is a
(faithfully) flat $R'$-module, then $M'$ is a
(faithfully) flat $R$-module.
\end{lemma}

\begin{proof}
The first statement of the lemma is a particular case of the
second, so it is clearly enough to prove the latter. Let
$R \to R'$ be a flat ring map, and $M'$ a flat $R'$-module.
We need to prove that $M'$ is a flat $R$-module. Let
$N_1 \to N_2 \to N_3$ be an exact complex of $R$-modules.
Then, the complex $R' \otimes_R N_1 \to
R' \otimes_R N_2 \to R' \otimes_R N_3$ is exact (since $R'$
is flat as an $R$-module), and so the complex
$M' \otimes_{R'} \left(R' \otimes_R N_1\right)
\to M' \otimes_{R'} \left(R' \otimes_R N_2\right)
\to M' \otimes_{R'} \left(R' \otimes_R N_3\right)$ is
exact (since $M'$ is a flat $R'$-module). Since
$M' \otimes_{R'} \left(R' \otimes_R N\right)
\cong \left(M' \otimes_{R'} R'\right) \otimes_R N
\cong M' \otimes_R N$ for any $R$-module $N$ functorially
(by Lemmas \ref{lemma-tensor-with-bimodule} and
\ref{lemma-flip-tensor-product}), this complex is isomorphic
to the complex
$M' \otimes_R N_1 \to M' \otimes_R N_2 \to M' \otimes_R N_3$,
which is therefore also exact. This shows that $M'$ is a flat
$R$-module. Tracing this argument backwards, we can show
that if $R \to R'$ is faithfully flat, and if $M'$ is
faithfully flat as an $R'$-module, then $M'$ is faithfully
flat as an $R$-module.
\end{proof}

\begin{lemma}
\label{lemma-flat}
Let $M$ be an $R$-module. The following are equivalent:
\begin{enumerate}
\item
\label{item-flat}
$M$ is flat over $R$.
\item
\label{item-injective}
for every injection of $R$-modules $N \subset N'$
the map $N \otimes_R M \to N'\otimes_R M$ is injective.
\item
\label{item-f-ideal}
for every ideal $I \subset R$ the map
$I \otimes_R M \to R \otimes_R M = M$ is injective.
\item
\label{item-ffg-ideal}
for every finitely generated ideal $I \subset R$
the map $I \otimes_R M \to R \otimes_R M = M$ is injective.
\end{enumerate}
\end{lemma}

\begin{proof}
The implications (\ref{item-flat}) implies (\ref{item-injective})
implies (\ref{item-f-ideal}) implies (\ref{item-ffg-ideal}) are all
trivial. Thus we prove (\ref{item-ffg-ideal}) implies (\ref{item-flat}).
Suppose that $N_1 \to N_2 \to N_3$ is exact.
Let $K = \Ker(N_2 \to N_3)$ and $Q = \Im(N_2 \to N_3)$.
Then we get maps
$$
N_1 \otimes_R M \to
K \otimes_R M \to
N_2 \otimes_R M \to
Q \otimes_R M \to
N_3 \otimes_R M
$$
Observe that the first and third arrows are surjective. Thus if we show
that the second and fourth arrows are injective, then we are
done\footnote{Here is the argument in more detail:
Assume that we know that the second and fourth arrows are
injective. Lemma \ref{lemma-tensor-product-exact} (applied
to the exact sequence $K \to N_2 \to Q \to 0$) yields that
the sequence $K \otimes_R M \to N_2 \otimes_R M \to
Q \otimes_R M \to 0$ is exact. Hence,
$\Ker \left(N_2 \otimes_R M \to Q \otimes_R M\right)
= \Im \left(K \otimes_R M \to N_2 \otimes_R M\right)$.
Since
$\Im \left(K \otimes_R M \to N_2 \otimes_R M\right)
= \Im \left(N_1 \otimes_R M \to N_2 \otimes_R M\right)$
(due to the surjectivity of $N_1 \otimes_R M \to
K \otimes_R M$) and
$\Ker \left(N_2 \otimes_R M \to Q \otimes_R M\right)
= \Ker \left(N_2 \otimes_R M \to N_3 \otimes_R M\right)$
(due to the injectivity of $Q \otimes_R M \to
N_3 \otimes_R M$), this becomes
$\Ker \left(N_2 \otimes_R M \to N_3 \otimes_R M\right)
= \Im \left(N_1 \otimes_R M \to N_2 \otimes_R M\right)$,
which shows that the functor $- \otimes_R M$ is exact,
whence $M$ is flat.}.
Hence it suffices to show that $- \otimes_R M$ transforms
injective $R$-module maps into injective $R$-module maps.

\medskip\noindent
Assume $K \to N$ is an injective $R$-module map and
let $x \in \Ker(K \otimes_R M \to N \otimes_R M)$.
We have to show that $x$ is zero.
The $R$-module $K$ is the union of its finite
$R$-submodules; hence, $K \otimes_R M$ is
the colimit of $R$-modules of the form
$K_i \otimes_R M$ where $K_i$ runs over all finite
$R$-submodules of $K$
(because tensor product commutes with colimits).
Thus, for some $i$ our $x$ comes from an element
$x_i \in K_i \otimes_R M$. Thus we may assume that $K$
is a finite $R$-module. Assume this. We regard the
injection $K \to N$ as an inclusion, so that
$K \subset N$.

\medskip\noindent
The $R$-module $N$ is the union of its finite
$R$-submodules that contain $K$. Hence, $N \otimes_R M$
is the colimit of $R$-modules of the form
$N_i \otimes_R M$ where $N_i$ runs over all finite
$R$-submodules of $N$ that contain $K$
(again since tensor product commutes with colimits).
Notice that this is a colimit over a directed system
(since the sum of two finite submodules of $N$ is
again finite).
Hence, (by Lemma \ref{lemma-zero-directed-limit})
the element $x \in K \otimes_R M$ maps to
zero in at least one of these $R$-modules
$N_i \otimes_R M$ (since $x$ maps to zero
in $N \otimes_R M$).
Thus we may assume $N$ is a finite $R$-module.

\medskip\noindent
Assume $N$ is a finite $R$-module. Write $N = R^{\oplus n}/L$ and $K = L'/L$
for some $L \subset L' \subset R^{\oplus n}$.
For any $R$-submodule $G \subset R^{\oplus n}$,
we have a canonical map $G \otimes_R M \to M^{\oplus n}$
obtained by composing
$G \otimes_R M \to R^n \otimes_R M = M^{\oplus n}$.
It suffices to prove that $L \otimes_R M \to M^{\oplus n}$
and $L' \otimes_R M \to M^{\oplus n}$ are injective.
Namely, if so, then we see that
$K \otimes_R M = L' \otimes_R M/L \otimes_R M \to M^{\oplus n}/L \otimes_R M$
is injective too\footnote{This becomes obvious if we
identify $L' \otimes_R M$ and $L \otimes_R M$ with
submodules of $M^{\oplus n}$ (which is legitimate since
the maps $L \otimes_R M \to M^{\oplus n}$
and $L' \otimes_R M \to M^{\oplus n}$ are injective and
commute with the obvious map $L' \otimes_R M \to L \otimes_R M$).}.

\medskip\noindent
Thus it suffices to show that $L \otimes_R M \to M^{\oplus n}$
is injective when $L \subset R^{\oplus n}$ is an $R$-submodule.
We do this by induction on $n$. The base case $n = 1$ we handle below.
For the induction step assume $n > 1$ and set
$L' = L \cap R \oplus 0^{\oplus n - 1}$. Then $L'' = L/L'$ is a submodule
of $R^{\oplus n - 1}$. We obtain a diagram
$$
\xymatrix{
&
L' \otimes_R M \ar[r] \ar[d] &
L \otimes_R M \ar[r] \ar[d] &
L'' \otimes_R M \ar[r] \ar[d] &
0 \\
0 \ar[r] &
M \ar[r] &
M^{\oplus n} \ar[r] &
M^{\oplus n - 1} \ar[r] & 0
}
$$
By induction hypothesis and the base case the left and right vertical
arrows are injective. The rows are exact. It follows that the middle vertical
arrow is injective too.

\medskip\noindent
The base case of the induction above is when $L \subset R$ is an ideal.
In other words, we have to show that $I \otimes_R M \to M$ is injective
for any ideal $I$ of $R$. We know this is true when $I$ is finitely
generated. However, $I = \bigcup I_\alpha$ is the union of the
finitely generated ideals $I_\alpha$ contained in it. In other words,
$I = \colim I_\alpha$. Since $\otimes$ commutes with colimits we see that
$I \otimes_R M = \colim I_\alpha \otimes_R M$ and since all
the morphisms $I_\alpha \otimes_R M \to M$ are injective by
assumption, the same is true for $I \otimes_R M \to M$.
\end{proof}

\begin{lemma}
\label{lemma-colimit-rings-flat}
Let $\{R_i, \varphi_{ii'}\}$ be a system of rings over the directed set $I$.
Let $R = \colim_i R_i$.
\begin{enumerate}
\item If $M$ is an $R$-module such that $M$ is flat as an $R_i$-module
for all $i$, then $M$ is flat as an $R$-module.
\item For $i \in I$ let $M_i$ be a flat $R_i$-module and
for $i' \geq i$ let $f_{ii'} : M_i \to M_{i'}$ be a $\varphi_{ii'}$-linear
map such that $f_{i' i''} \circ f_{i i'} = f_{i i''}$. Then
$M = \colim_{i \in I} M_i$ is a flat $R$-module.
\end{enumerate}
\end{lemma}

\begin{proof}
Part (1) is a special case of part (2) with $M_i = M$ for all $i$
and $f_{i i'} = \text{id}_M$. Proof of (2).
Let $\mathfrak a \subset R$ be a finitely generated ideal. By
Lemma \ref{lemma-flat}
it suffices to show that $\mathfrak a \otimes_R M \to M$ is
injective. We can find an $i \in I$ and a finitely generated ideal
$\mathfrak a' \subset R_i$ such that $\mathfrak a = \mathfrak a'R$.
Then $\mathfrak a = \colim_{i' \geq i} \mathfrak a'R_{i'}$.
Since $\otimes$ commutes with colimits the map $\mathfrak a \otimes_R M \to M$
is the colimit of the maps
$$
\mathfrak a'R_{i'} \otimes_{R_{i'}} M_{i'} \longrightarrow M_{i'}
$$
These maps are all injective by assumption. Since colimits over $I$
are exact by Lemma \ref{lemma-directed-colimit-exact} we win.
\end{proof}

\begin{lemma}
\label{lemma-flat-base-change}
Suppose that $M$ is (faithfully) flat over $R$, and that $R \to R'$
is a ring map. Then $M \otimes_R R'$ is (faithfully) flat over $R'$.
\end{lemma}

\begin{proof}
For any $R'$-module $N$ we have a canonical
isomorphism $N \otimes_{R'} (R'\otimes_R M)
= N \otimes_R M$. Hence the desired exactness properties of the functor
$-\otimes_{R'}(R'\otimes_R M)$ follow from
the corresponding exactness properties of the functor $-\otimes_R M$.
\end{proof}

\begin{lemma}
\label{lemma-flatness-descends}
Let $R \to R'$ be a faithfully flat ring map.
Let $M$ be a module over $R$, and set $M' = R' \otimes_R M$.
Then $M$ is flat over $R$ if and only if $M'$ is flat over $R'$.
\end{lemma}

\begin{proof}
By Lemma \ref{lemma-flat-base-change} we see that if $M$ is flat
then $M'$ is flat. For the converse, suppose that $M'$ is flat.
Let $N_1 \to N_2 \to N_3$ be an exact sequence of $R$-modules.
We want to show that $N_1 \otimes_R M \to N_2 \otimes_R M \to N_3 \otimes_R M$
is exact. We know that
$N_1 \otimes_R R' \to N_2 \otimes_R R' \to N_3 \otimes_R R'$ is
exact, because $R \to R'$ is flat. Flatness of $M'$ implies that
$N_1 \otimes_R R' \otimes_{R'} M'
\to N_2 \otimes_R R' \otimes_{R'} M'
\to N_3 \otimes_R R' \otimes_{R'} M'$ is exact.
We may write this as
$N_1 \otimes_R M \otimes_R R'
\to N_2 \otimes_R M \otimes_R R'
\to N_3 \otimes_R M \otimes_R R'$.
Finally, faithful flatness implies that
$N_1 \otimes_R M \to N_2 \otimes_R M \to N_3 \otimes_R M$
is exact.
\end{proof}

\begin{lemma}
\label{lemma-flatness-descends-more-general}
Let $R$ be a ring. Let $S \to S'$ be a flat map of $R$-algebras.
Let $M$ be a module over $S$, and set $M' = S' \otimes_S M$.
\begin{enumerate}
\item If $M$ is flat over $R$, then $M'$ is flat over $R$.
\item If $S \to S'$ is faithfully flat, then $M$ is flat
over $R$ if and only if $M'$ is flat over $R$.
\end{enumerate}
\end{lemma}

\begin{proof}
Let $N \to N'$ be an injection of $R$-modules. By the flatness
of $S \to S'$ we have
$$
\Ker(N \otimes_R M \to N' \otimes_R M) \otimes_S S'
=
\Ker(N \otimes_R M' \to N' \otimes_R M')
$$
If $M$ is flat over $R$, then the left hand side is zero and
we find that $M'$ is flat over $R$ by the second characterization
of flatness in Lemma \ref{lemma-flat}.
If $M'$ is flat over $R$ then we have the vanishing of the right hand side
and if in addition $S \to S'$ is faithfully flat, this implies that
$\Ker(N \otimes_R M \to N' \otimes_R M)$ is zero which in turn
shows that $M$ is flat over $R$.
\end{proof}

\begin{lemma}
\label{lemma-flat-permanence}
Let $R \to S$ be a ring map. Let $M$ be an $S$-module.
If $M$ is flat as an $R$-module and faithfully flat as an $S$-module,
then $R \to S$ is flat.
\end{lemma}

\begin{proof}
Let $N_1 \to N_2 \to N_3$ be an exact sequence of $R$-modules.
By assumption $N_1 \otimes_R M \to N_2 \otimes_R M \to N_3 \otimes_R M$
is exact. We may write this as
$$
N_1 \otimes_R S \otimes_S M
\to
N_2 \otimes_R S \otimes_S M
\to
N_3 \otimes_R S \otimes_S M.
$$
By faithful flatness of $M$ over $S$ we conclude that
$N_1 \otimes_R S \to N_2 \otimes_R S \to N_3 \otimes_R S$ is exact.
Hence $R \to S$ is flat.
\end{proof}

\noindent
Let $R$ be a ring.
Let $M$ be an $R$-module.
Let $\sum f_i x_i = 0$ be a relation in $M$.
We say the relation $\sum f_i x_i$
is {\it trivial} if there exist an integer $m \geq 0$,
elements $y_j \in M$, $j = 1, \ldots, m$, and elements $a_{ij} \in R$,
$i = 1, \ldots, n$, $j = 1, \ldots, m$ such that
$$
x_i = \sum\nolimits_j a_{ij} y_j, \forall i,
\quad\text{and}\quad
0 = \sum\nolimits_i f_ia_{ij}, \forall j.
$$

\begin{lemma}[Equational criterion of flatness]
\label{lemma-flat-eq}
A module $M$ over $R$ is flat if and only if
every relation in $M$ is trivial.
\end{lemma}

\begin{proof}
Assume $M$ is flat and let $\sum f_i x_i = 0$ be a relation in $M$.
Let $I = (f_1, \ldots, f_n)$, and let
$K = \Ker(R^n \to I, (a_1, \ldots, a_n) \mapsto \sum_i a_i f_i)$.
So we have the short exact sequence
$0 \to K \to R^n \to I \to 0$. Then $\sum f_i \otimes x_i$
is an element of $I \otimes_R M$ which maps
to zero in $R \otimes_R M = M$. By flatness
$\sum f_i \otimes x_i$ is zero in $I \otimes_R M$.
Thus there exists an element of $K \otimes_R M$ mapping
to $\sum e_i \otimes x_i \in R^n \otimes_R M$ where $e_i$
is the $i$th basis element of $R^n$.
Write this element as $\sum k_j \otimes y_j$
and then write the image of $k_j$ in $R^n$ as
$\sum a_{ij} e_i$ to get the result.

\medskip\noindent
Assume every relation is trivial, let $I$
be a finitely generated ideal, and let $x = \sum f_i \otimes x_i$
be an element of $I \otimes_R M$ mapping to zero in $R \otimes_R M = M$.
This just means exactly that $\sum f_i x_i$ is a relation in
$M$. And the fact that it is trivial implies easily that
$x$ is zero, because
$$
x
=
\sum f_i \otimes x_i
=
\sum f_i \otimes \left(\sum a_{ij}y_j\right)
=
\sum \left(\sum f_i a_{ij}\right) \otimes y_j
=
0
$$
\end{proof}

\begin{lemma}
\label{lemma-flat-tor-zero}
Suppose that $R$ is a ring, $0 \to M'' \to M' \to M \to 0$
a short exact sequence, and $N$ an $R$-module. If $M$ is flat
then $N \otimes_R M'' \to N \otimes_R M'$ is injective, i.e., the
sequence
$$
0 \to N \otimes_R M'' \to N \otimes_R M' \to N \otimes_R M \to 0
$$
is a short exact sequence.
\end{lemma}

\begin{proof}
Let $R^{(I)} \to N$ be a surjection from a free module
onto $N$ with kernel $K$. The result follows
from the snake lemma applied to the following diagram
$$
\begin{matrix}
 & & 0 & & 0 & & 0 & & \\
 & & \uparrow & & \uparrow & & \uparrow & & \\
 & & M''\otimes_R N & \to & M' \otimes_R N & \to & M \otimes_R N & \to & 0 \\
 & & \uparrow & & \uparrow & & \uparrow & & \\
0 & \to & (M'')^{(I)} & \to & (M')^{(I)} & \to & M^{(I)} & \to & 0 \\
 & & \uparrow & & \uparrow & & \uparrow & & \\
 & & M''\otimes_R K & \to & M' \otimes_R K & \to & M \otimes_R K & \to & 0 \\
 & & & & & & \uparrow & & \\
 & & & & & & 0 & &
\end{matrix}
$$
with exact rows and columns. The middle row is exact because tensoring
with the free module $R^{(I)}$ is exact.
\end{proof}


\begin{lemma}
\label{lemma-flat-ses}
Suppose that $0 \to M' \to M \to M'' \to 0$ is
a short exact sequence of $R$-modules.
If $M'$ and $M''$ are flat so is $M$.
If $M$ and $M''$ are flat so is $M'$.
\end{lemma}

\begin{proof}
We will use the criterion that a module $N$ is flat if for
every ideal $I \subset R$ the map $N \otimes_R I \to N$ is injective,
see Lemma \ref{lemma-flat}.
Consider an ideal $I \subset R$.
Consider the diagram
$$
\begin{matrix}
0 & \to & M' & \to & M & \to & M'' & \to & 0 \\
& & \uparrow & & \uparrow & & \uparrow & & \\
& & M'\otimes_R I & \to & M \otimes_R I & \to & M''\otimes_R I & \to & 0
\end{matrix}
$$
with exact rows. This immediately proves the first assertion.
The second follows because if $M''$ is flat then the lower left
horizontal arrow is injective by Lemma \ref{lemma-flat-tor-zero}.
\end{proof}

\begin{lemma}
\label{lemma-easy-ff}
Let $R$ be a ring.
Let $M$ be an $R$-module.
The following are equivalent
\begin{enumerate}
\item $M$ is faithfully flat, and
\item $M$ is flat and for all $R$-module homomorphisms $\alpha : N \to N'$
we have $\alpha = 0$ if and only if $\alpha \otimes \text{id}_M = 0$.
\end{enumerate}
\end{lemma}

\begin{proof}
If $M$ is faithfully flat, then
$0 \to \Ker(\alpha) \to N \to N'$ is exact if and only if the same holds
after tensoring with $M$. This proves (1) implies (2).
For the other, assume (2). Let $N_1 \to N_2 \to N_3$
be a complex, and assume the complex
$N_1 \otimes_R M \to N_2 \otimes_R M \to N_3\otimes_R M$
is exact. Take $x \in \Ker(N_2 \to N_3)$,
and consider the map $\alpha : R \to N_2/\Im(N_1)$,
$r \mapsto rx + \Im(N_1)$. By the exactness
of the complex $-\otimes_R M$ we see that $\alpha \otimes
\text{id}_M$ is zero. By assumption we get that $\alpha$ is
zero. Hence $x $ is in the image of $N_1 \to N_2$.
\end{proof}

\begin{lemma}
\label{lemma-ff}
\begin{slogan}
A flat module is faithfully flat if and only if it has nonzero fibers.
\end{slogan}
Let $M$ be a flat $R$-module.
The following are equivalent:
\begin{enumerate}
\item $M$ is faithfully flat,
\item for every nonzero $R$-module $N$, then tensor product $M \otimes_R N$
is nonzero,
\item for all $\mathfrak p \in \Spec(R)$
the tensor product $M \otimes_R \kappa(\mathfrak p)$ is nonzero, and
\item for all maximal ideals $\mathfrak m$ of $R$
the tensor product $M \otimes_R \kappa(\mathfrak m) = M/{\mathfrak m}M$
is nonzero.
\end{enumerate}
\end{lemma}

\begin{proof}
Assume $M$ faithfully flat and $N \not = 0$. By Lemma \ref{lemma-easy-ff}
the nonzero map $1 : N \to N$ induces a nonzero map
$M \otimes_R N \to M \otimes_R N$, so $M \otimes_R N \not = 0$.
Thus (1) implies (2). The implications (2) $\Rightarrow$ (3) $\Rightarrow$ (4)
are immediate.

\medskip\noindent
Assume (4). Suppose that $N_1 \to N_2 \to N_3$ is a complex and
suppose that $N_1 \otimes_R M \to N_2\otimes_R M \to
N_3\otimes_R M$ is exact. Let $H$ be the cohomology of the complex,
so $H = \Ker(N_2 \to N_3)/\Im(N_1 \to N_2)$. To finish the proof
we will show $H = 0$. By flatness we see that $H \otimes_R M = 0$.
Take $x \in H$ and let $I = \{f \in R \mid fx = 0 \}$
be its annihilator. Since $R/I \subset H$ we get
$M/IM \subset H \otimes_R M = 0$ by flatness of $M$.
If $I \not =  R$ we may choose
a maximal ideal $I \subset \mathfrak m \subset R$.
This immediately gives a contradiction.
\end{proof}

\begin{lemma}
\label{lemma-ff-rings}
Let $R \to S$ be a flat ring map.
The following are equivalent:
\begin{enumerate}
\item $R \to S$ is faithfully flat,
\item the induced map on $\Spec$ is surjective, and
\item any closed point $x \in \Spec(R)$ is
in the image of the map $\Spec(S) \to \Spec(R)$.
\end{enumerate}
\end{lemma}

\begin{proof}
This follows quickly from Lemma \ref{lemma-ff}, because we
saw in Remark \ref{remark-fundamental-diagram}
that $\mathfrak p$ is in the image
if and only if the ring $S \otimes_R \kappa(\mathfrak p)$
is nonzero.
\end{proof}

\begin{lemma}
\label{lemma-local-flat-ff}
A flat local ring homomorphism of local rings is faithfully flat.
\end{lemma}

\begin{proof}
Immediate from Lemma \ref{lemma-ff-rings}.
\end{proof}

\noindent
Flatness meshes well with localization.

\begin{lemma}
\label{lemma-flat-localization}
Let $R$ be a ring. Let $S \subset R$ be a multiplicative subset.
\begin{enumerate}
\item The localization $S^{-1}R$ is a flat $R$-algebra.
\item If $M$ is an $S^{-1}R$-module, then $M$ is a flat $R$-module
if and only if $M$ is a flat $S^{-1}R$-module.
\item Suppose $M$ is an $R$-module. Then
$M$ is a flat $R$-module if and only if $M_{\mathfrak p}$ is a flat
$R_{\mathfrak p}$-module for all primes $\mathfrak p$ of $R$.
\item Suppose $M$ is an $R$-module. Then $M$ is a flat $R$-module if
and only if $M_{\mathfrak m}$ is a flat
$R_{\mathfrak m}$-module for all maximal ideals $\mathfrak m$ of $R$.
\item Suppose $R \to A$ is a ring map, $M$ is an $A$-module,
and $g_1, \ldots, g_m \in A$ are elements generating the unit
ideal of $A$. Then $M$ is flat over $R$ if and only if each localization
$M_{g_i}$ is flat over $R$.
\item Suppose $R \to A$ is a ring map, and $M$ is an $A$-module.
Then $M$ is a flat $R$-module if and only if the localization
$M_{\mathfrak q}$ is a flat $R_{\mathfrak p}$-module
(with $\mathfrak p$ the prime of $R$ lying under $\mathfrak q$)
for all primes $\mathfrak q$ of $A$.
\item Suppose $R \to A$ is a ring map, and $M$ is an $A$-module.
Then $M$ is a flat $R$-module if and only if the localization
$M_{\mathfrak m}$ is a flat $R_{\mathfrak p}$-module
(with $\mathfrak p = R \cap \mathfrak m$)
for all maximal ideals $\mathfrak m$ of $A$.
\end{enumerate}
\end{lemma}

\begin{proof}
Let us prove the last statement of the lemma.
In the proof we will use repeatedly that localization is exact
and commutes with tensor product, see Sections \ref{section-localization}
and \ref{section-tensor-product}.

\medskip\noindent
Suppose $R \to A$ is a ring map, and $M$ is an $A$-module.
Assume that $M_{\mathfrak m}$ is a flat $R_{\mathfrak p}$-module
for all maximal ideals $\mathfrak m$ of $A$ (with
$\mathfrak p = R \cap \mathfrak m$). Let $I \subset R$ be an ideal.
We have to show the map $I \otimes_R M \to M$ is injective.
We can think of this as a map of $A$-modules.
By assumption the localization
$(I \otimes_R M)_{\mathfrak m} \to M_{\mathfrak m}$ is injective
because
$(I \otimes_R M)_{\mathfrak m} =
I_{\mathfrak p} \otimes_{R_{\mathfrak p}} M_{\mathfrak m}$.
Hence the kernel of $I \otimes_R M \to M$ is zero by
Lemma \ref{lemma-characterize-zero-local}.
Hence $M$ is flat over $R$.

\medskip\noindent
Conversely, assume $M$ is flat over $R$. Pick a prime $\mathfrak q$
of $A$ lying over the prime $\mathfrak p$ of $R$. Suppose that
$I \subset R_{\mathfrak p}$ is an ideal. We have to show that
$I \otimes_{R_{\mathfrak p}} M_{\mathfrak q} \to M_{\mathfrak q}$
is injective. We can write $I = J_{\mathfrak p}$ for some
ideal $J \subset R$. Then the map
$I \otimes_{R_{\mathfrak p}} M_{\mathfrak q} \to M_{\mathfrak q}$
is just the localization (at $\mathfrak q$) of the map
$J \otimes_R M \to M$ which is injective. Since localization is exact
we see that $M_{\mathfrak q}$ is a flat $R_{\mathfrak p}$-module.

\medskip\noindent
This proves (7) and (6). The other statements follow in a straightforward
way from the last statement (proofs omitted).
\end{proof}

\begin{lemma}
\label{lemma-flat-going-down}
Let $R \to S$ be flat. Let $\mathfrak p \subset \mathfrak p'$
be primes of $R$. Let $\mathfrak q' \subset S$ be a prime of $S$
mapping to $\mathfrak p'$. Then there exists a prime
$\mathfrak q \subset \mathfrak q'$ mapping to $\mathfrak p$.
\end{lemma}

\begin{proof}
By Lemma \ref{lemma-flat-localization} the local ring map
$R_{\mathfrak p'} \to S_{\mathfrak q'}$ is flat.
By Lemma \ref{lemma-local-flat-ff} this local ring map is faithfully
flat. By Lemma \ref{lemma-ff-rings} there is a prime mapping to
$\mathfrak p R_{\mathfrak p'}$. The inverse image of this
prime in $S$ does the job.
\end{proof}

\noindent
The property of $R \to S$ described in the lemma is called the
``going down property''. See Definition \ref{definition-going-up-down}.

\begin{lemma}
\label{lemma-colimit-faithfully-flat}
Let $R$ be a ring. Let $\{S_i, \varphi_{ii'}\}$ be a directed system of
faithfully flat $R$-algebras. Then $S = \colim_i S_i$ is a faithfully flat
$R$-algebra.
\end{lemma}

\begin{proof}
By Lemma \ref{lemma-colimit-flat} we see that $S$ is flat.
Let $\mathfrak m \subset R$ be a maximal ideal. By
Lemma \ref{lemma-ff-rings}
none of the rings $S_i/\mathfrak m S_i$ is zero.
Hence $S/\mathfrak mS = \colim S_i/\mathfrak mS_i$ is nonzero
as well because $1$ is not equal to zero. Thus the image of
$\Spec(S) \to \Spec(R)$ contains $\mathfrak m$ and we see that $R \to S$
is faithfully flat by Lemma \ref{lemma-ff-rings}.
\end{proof}





